\section{Related literature}\label{sec:lit}

\subsection{AI as an input to climate adaptation, and its costs}

Adaptation is an information problem as much as an investment problem. Households, firms and
governments must anticipate hazards, identify who is exposed and allocate scarce resources before
and after shocks. Machine learning lowers the cost of these tasks. It improves forecasts and early
warnings of compound events \citep{reichstein2025}, extracts crop, land-use and settlement
information from satellite imagery, and supports the targeting of relief
\citep{rolnick2022,hilbert2016}. Reviews of AI for the Sustainable Development Goals identify climate
action and food security among the goals AI most often enables \citep{vinuesa2020}, and reviews of
AI for adaptation document applications in agriculture, water, health and infrastructure
\citep{lealfilho2022}.

Domestic capability matters for these benefits. Models trained on data from other regions perform
poorly in data-scarce settings, and the agencies that run early-warning, extension and
disaster-risk-finance systems need local analytical skills to adapt, validate and maintain them
\citep{cowls2023,larosa2024}. AI is a general-purpose technology whose returns depend on
complementary investments and arrive with delay \citep{brynjolfsson2021}. A country's research
base in climate-relevant AI is therefore a plausible, if imperfect, indicator of its capacity to
turn AI into adaptation, and whether that research is led domestically or by foreign partners is a
plausible indicator of whether it addresses local priorities.

The costs are also documented. Training and running large models uses electricity and water
\citep{strubell2019,devries2023}. Data-centre demand is projected to rise steeply \citep{iea2025},
after a decade in which efficiency gains held it roughly flat \citep{masanet2020}. At the economy
level, digitalisation has ambiguous effects on energy demand \citep{lange2020}. Studies of Chinese
industries and cities find that AI lowers carbon intensity \citep{liu2022,chen2022}, cross-country
work finds mixed effects \citep{wang2024}, and \citet{dong2024} show that income inequality weakens
AI's abatement benefits. The system-level effect depends on what AI is used for \citep{kaack2022}.

\subsection{Who bears climate damages}

Climate econometrics uses within-country variation in temperature to estimate damages. Hot years
reduce growth in poor countries \citep{dell2012}, the global response of output to temperature is
non-linear \citep{burke2015}, and damages vary across regions with climate and income
\citep{kalkuhl2020,hsiang2016}. Warming has widened the gap between rich and poor countries
\citep{diffenbaugh2019}, disasters kill more people where incomes are low and institutions weak
\citep{kahn2005}, and their macroeconomic costs are larger in poorer economies \citep{noy2009}.
Within countries, poor households are more exposed, lose a larger share of their assets and recover
more slowly \citep{hallegatte2017}. Agriculture is the channel through which heat most directly
reaches the poor. Crop yields fall sharply once temperatures exceed crop-specific thresholds
\citep{schlenker2009}, warming since 1980 has already reduced global maize and wheat production
\citep{lobell2011}, and anthropogenic climate change has slowed agricultural productivity growth,
most of all in warmer, poorer regions \citep{ortizbobea2021}. Because the poorest households in
developing economies depend disproportionately on farming, and agricultural growth reduces poverty
more than growth in other sectors \citep{christiaensen2011}, heat damage to yields is a damage to
the bottom of the distribution. Little of this work asks which capabilities reduce damages, and
almost none of it measures damages below the country level for a global sample.

The literature on AI and inequality suggests why AI capability might protect some groups more than
others. Automation and AI raise returns to capital and to skills that are unevenly held
\citep{acemoglu2018,acemoglu2020,felten2021}. About 40 per cent of jobs worldwide are exposed to AI,
with complementarity concentrated among high-income workers \citep{cazzaniga2024}, and AI may widen
the gap between advanced and developing economies \citep{korinek2021}. The value of AI-based
warnings and advice depends on connectivity, literacy and trust \citep{worldbank2016}. Where these
are unequally distributed, the benefits of AI-enabled adaptation may accrue to those who were
already less vulnerable \citep{larosa2024}.

\subsection{From heterogeneous effects to equity-minded targeting}

The empirical studies cited above mostly estimate an average effect and test for heterogeneity with
an interaction term, for example between AI and the Gini coefficient \citep{dong2024}. That approach
imposes the dimension of heterogeneity in advance and cannot tell whether AI matters or merely
proxies for income or institutions. It also answers a different question from the one a planner
faces. Recent econometric work offers better tools. Cross-fitted, Neyman-orthogonal scores allow
flexible machine-learning estimation of nuisance functions without biasing low-dimensional targets
\citep{robinson1988,chernozhukov2018}. Causal forests and their generalisations estimate
conditional effects \citep{wager2018,athey2019}, and generic machine-learning inference provides
valid statements about which groups benefit most \citep{chernozhukov2025}.

Treatment-choice methods go one step further and estimate the rule itself. Building on
\citet{manski2004}, empirical welfare maximisation \citep{kitagawa2018} and doubly robust policy
learning \citep{athey2021} choose an assignment rule from a constrained class with guaranteed
regret. \citet{kitagawa2021} extend the approach to rank-dependent social welfare functions that
weight outcomes by their position in the distribution. These methods are designed for a treatment
whose effect on an outcome is identified directly. In our setting the ``treatment'', climate-AI
capability, is not randomly assigned. What is identified is how it changes the effect of a
different, exogenous variable, the temperature shock, and the welfare-relevant outcome is observed
separately for each income decile. EW-DRPL adapts orthogonal scores and equity-weighted policy
learning to this moderated, decile-level structure.
